% page 270 of the facsimile (image p-269 of the scan)
% block 157.5 x 250.6 mm ; left margin 8.0 mm
\pgc{-12.700mm}{0.000mm}{
\l{\cel{5}262}
\l{}
\l{\sou{kapitulo}. I. (\sou{bot}.) Infloresenco sama-forma kam kapo, spiko.}
\l{\cel{3}- II. (\sou{liturgio katol}.) Kurta fragmento di la santa skri-}
\l{\cel{3}buri,qua relatas la oficio di la dio, quan onu recitas}
\l{\cel{3}pos la psalmi, la precipua kloki di la breviario.-DEF.}
\l{}
\l{\sou{kapitulacar}.\cel{2}(\sou{netrans.})\cel{2}Kontratar kun la enemiko, pri la}
\l{\cel{3}kondicioni di la livro (forme di rezisto-ceso) di forti-}
\l{\cel{3}fikajo, fortreso, armeo. - DEFIRS.}
\l{}
\l{\sou{kapono}. Hanulo yuna quan onu kastris por igar lu divenar}
\l{\cel{3}plu grasoza. - DEFIRS.}
\l{}
\l{\sou{kaponiero}. (\sou{milit-arto}). Galerio konstruktita en sika fosa-}
\l{\cel{3}to di fortreso, infre di la pafo-lineo, por komunikar}
\l{\cel{3}de ta a ca fortifikajo, e por impedar la transiro di la}
\l{\cel{3}fosato. - DEFIRS.}
\l{}
\l{\sou{kaporalo}. Armeano di qua la grado esas la minim alta, -}
\l{\cel{2}DEFIR.}
\l{}
\l{\sou{kapoto}. I. Longa redingoto por la viri, e precipue por la}
\l{\cel{3}soldati. - II. l. sur ula komerco-navo : bloko qua shir-}
\l{\cel{3}mas la enireyo di la pupo-eskalero ; 2. sur la automobi-}
\l{\cel{3}li, kovrilo metala qua shirmas la motoro. - DEFIRS.}
\l{}
\l{\sou{kapro}. (\sou{zool.})\cel{2}Mamifero, speco ek la familio \textquotedbl{}rumineri\textquotedbl{},}
\l{\cel{3}kun korni kava e vivaca, mentono kun barbo, e di qua la}
\l{\cel{3}maskulo odoras multe e desagreable. - L.\cel{1}\sou{capra (hircus})}
\l{\cel{3}- FIS.}
\l{}
\l{\sou{kapreolo}. (\sou{zool.}) Speco di cervo, ma plu mikra, di qua la}
\l{\cel{3}korni esas kurta, cilindra, ed havas nur un brancho. -}
\l{\cel{3}L.\cel{1}\sou{cervus capreolus}. - IL.}
\l{}
\l{\sou{kaprico}. Volo subita, variema, nejustifikata, fantazia.-DEFIRS}
\l{}
\l{\sou{kaprifolio}. (\sou{bot.)}\cel{2}Arboreto klimera, sarmentoza, di qua la}
\l{\cel{3}flori odoras. - L.\cel{1}\sou{lonicera caprifolium}. - DIR.}
\l{}
\l{\sou{kaprikorno}. I. (\sou{zool.}) Insekto koleoptera, kun longa anteni}
\l{\cel{3}artikoza, kurva quale korni. - L.\cel{1}\sou{cerembyx}. - II. Stelaro}
\l{\cel{3}imaginata kom kaproforma, qua okupas ta parto di zodiako}
\l{\cel{3}quan enirasSuno lor la solstico vintrala, e quan lu sem-}
\l{\cel{3}blas transirar inter la 20esma di decembro e la 20esma}
\l{\cel{3}di januaro. - EFIS.}
\l{}
\l{\sou{kaprimulgo}. (\sou{zool.})\cel{2}Genero di pasero kun beko fendita, e}
\l{\cel{3}qua, flugante, havas sua beko apertita. - L.\cel{1}\sou{caprimalgus}}
\l{\cel{3}\sou{europeus}. - IL.}
\l{}
\l{\sou{kapriolar}. (\sou{netrans.}) I. Saltar bruske e retrosaltar, lor}
\l{\cel{3}folo-ludeto. - II. (\sou{aludante kavalo)}. Saltar e, dum esar}
\l{\cel{3}super sulo, kikar. - DEFIRS.}
\l{}
\l{\sou{kapstano}.\cel{2}Vindilo vertikala, an-cirkum qua spulesas kablo.}
\l{\cel{3}DERS.}
}
